% ==== 0 ====
\yoschapter{0}{Preliminaries}
\label{chap:0}

It is the purpose of this chapter to explain certain notions and theorems used throughout the present book. These are related to Set Theory, Topological Spaces, Measure Spaces and Linear Spaces. 


\yossection{1}{Set Theory}
\label{sec:0.1}

Sets. $x \in X$ means that $x$ is a member or element of the set $X; x \notin X$ means that $x$ is not a member of the set $X$. We denote the set consisting of all $x$ possessing the property $P$ by $\{x; P\}$. Thus $\{y; y = x\}$ is the set $\{x\}$ consisting of a single element $x$. The void set is the set with no members, and will be denoted by $\emptyset$. If every element of a set $X$ is also an element of a set $Y$, then $X$ is said to be a subset of $Y$ and this fact will be denoted by $X \subseteq Y$, or $Y \supseteq X$. If $\mathfrak{X}$ is a set whose elements are sets $X$, then the set of all $x$ such that $x \in X$ for some $X \in \mathfrak{X}$ is called the union of sets $X$ in $\mathfrak{X}$; this union will be denoted by $\bigcup\limits_{X \in \mathfrak{X}} X$. The intersection of the sets $X$ in $\mathfrak{X}$ is the set of all $x$ which are elements of every $X \in \mathfrak{X}$; this intersection will be denoted by $\bigcap\limits_{X \in \mathfrak{X}} X$. Two sets are disjoint if their intersection is void. A family of sets is disjoint if every pair of distinct sets in the family is disjoint. If a sequence $\{X_n\}_{n=1,2,\ldots}$ of sets is a disjoint family, then the union $\bigcup\limits_{n=1}^{\infty} X_n$ may be written in the form of a sum $\sum\limits_{n=1}^{\infty} X_n$. 

Mappings. The term mapping, function and transformation will be used synonymously. The symbol $f: X \to Y$ will mean that $f$ is a single-valued function whose domain is $X$ and whose range is contained in $Y$; for every $x \in X$, the function $f$ assigns a uniquely determined element $f(x) = y \in Y$. For two mappings $f: X \to Y$ and $g: Y \to Z$, we can define their composite mapping $gf: X \to Z$ by $(gf)(x) = g(f(x))$. The symbol $f(M)$ denotes the set $\{f(x); x \in M\}$ and $f(M)$ is called the image of $M$ under the mapping $f$. The symbol $f^{-1}(N)$ denotes the set $\{x; f(x) \in N\}$ and $f^{-1}(N)$ is called the inverse image of $N$ under the mapping $f$. It is clear that 
$Y_{1} = f(f^{-1}(Y_{1}))$ for all $Y_{1} \subseteq f(X)$, and $X_{1} \subseteq f^{-1}(f(X_{1}))$ for all $X_{1} \subseteq X$.

If $f: X \to Y$, and for each $y \in f(X)$ there is only one $x \in X$ with $f(x) = y$, then $f$ is said to have an inverse (mapping) or to be one-to-one. The inverse mapping then has the domain $f(X)$ and range $X$; it is defined by the equation $x = f^{-1}(y) = f^{-1}(\{y\})$. 

The domain and the range of a mapping $f$ will be denoted by $D(f)$ and $R(f)$, respectively. Thus, if $f$ has an inverse then 
\[
f ^{-1} (f (x)) = x \text{ for all } x \in D (f), \text{ and } f (f ^{-1} (y)) = y \text{ for all } y \in R (f).
\]

The function $f$ is said to map $X$ onto $Y$ if $f(X) = Y$ and into $Y$ if $f(X) \subseteq Y$. The function $f$ is said to be an extension of the function $g$ and $g$ a restriction of $f$ if $D(f)$ contains $D(g)$, and $f(x) = g(x)$ for all $x$ in $D(g)$. 

\yossub{Zorn's Lemma}

\begin{definitionu} Let $P$ be a set of elements $a, b, \ldots$. Suppose there is a binary relation defined between certain pairs $(a, b)$ of elements of $P$, expressed by $a \prec b$, with the properties: 
$a\prec a$; if $a\prec b$ and $b\prec a$, then $a=b$; if $a\prec b$ and $b\prec c$, then $a\prec c$ (transitivity).
\end{definitionu}

Then $P$ is said to be partially ordered (or semi-ordered) by the relation $\prec$. 

\begin{examplen} If $P$ is the set of all subsets of a given set $X$, then the set inclusion relation $(A \subseteq B)$ gives a partial ordering of $P$. The set of all complex numbers $z = x + iy$, $w = u + iv$, $\ldots$ is partially ordered by defining $z \prec w$ to mean $x \leqq u$ and $y \leqq v$.

\end{examplen}
\begin{definitionu} Let $P$ be a partially ordered set with elements $a, b, \ldots$. If $a < c$ and $b < c$, we call $c$ an upper bound for $a$ and $b$. If furthermore $c < d$ whenever $d$ is an upper bound for $a$ and $b$, we call $c$ the least upper bound or the supremum of $a$ and $b$, and write $c = \sup (a, b)$ or $a \vee b$. This element of $P$ is unique if it exists. In a similar way we define the greatest lower bound or the infimum of $a$ and $b$, and denote it by $\inf (a, b)$ or $a \wedge b$. If $a \vee b$ and $a \wedge b$ exist for every pair $(a, b)$ in a partially ordered set $P$, $P$ is called a lattice. 
\end{definitionu}

\begin{examplen} The totality of subsets $M$ of a fixed set $B$ is a lattice by the partial ordering $M_{1} \prec M_{2}$ defined by the set inclusion relation $M_{1} \subseteq M_{2}$.

\end{examplen}
\begin{definitionu} $A$ partially ordered set $P$ is said to be linearly ordered (or totally ordered) if for every pair $(a, b)$ in $P$, either $a < b$ or $b < a$ holds. A subset of a partially ordered set is itself partially ordered by the relation which partially orders $P$; the subset might turn out to be linearly ordered by this relation. If $P$ is partially ordered and $S$ is a subset of $P$, an $m \in P$ is called an upper bound of $S$ if $s < m$ for every $s \in S$. An $m \in P$ is said to be maximal if $p \in P$ and $m < p$ together imply $m = p$. 
\end{definitionu}

\begin{lemman}[Zorn's Lemma] Let $P$ be a non-empty partially ordered set with the property that every linearly ordered subset of $P$ has an upper bound in $P$. Then $P$ contains at least one maximal element. 
\end{lemman}

It is known that Zorn's lemma is equivalent to Zermelo's axiom of choice in set theory. 


\yossection{2}{Topological Spaces}
\label{sec:0.2}

\yossub{Open Sets and Closed Sets}

\begin{definitionu} $A$ system $\tau$ of subsets of a set $X$ defines a topology in $X$ if $\tau$ contains the void set, the set $X$ itself, the union of every one of its subsystems, and the intersection of every one of its finite subsystems. The sets in $\tau$ are called the open sets of the topological space $(X,\tau)$; we shall often omit $\tau$ and refer to $X$ as a topological space. Unless otherwise stated, we shall assume that a topological space $X$ satisfies Hausdorff's axiom of separation: 

For every pair $(x_{1}, x_{2})$ of distinct points $x_{1}, x_{2}$ of $X$, there exist disjoint open sets $G_{1}, G_{2}$ such that $x_{1} \in G_{1}, x_{2} \in G_{2}$. 
\end{definitionu}

A neighbourhood of the point $x$ of $X$ is a set containing an open set which contains $x$. A neighbourhood of the subset $M$ of $X$ is a set which is a neighbourhood of every point of $M$. A point $x$ of $X$ is an accumulation point or limit point of a subset $M$ of $X$ if every neighbourhood of $x$ contains at least one point $m \in M$ different from $x$. 

\begin{definitionu} Any subset $M$ of a topological space $X$ becomes a topological space by calling ``open'' the subsets of $M$ which are of the form $M \cap G$ where $G'$'s are open sets of $X$. The induced topology of $M$ is called the relative topology of $M$ as a subset of the topological space $X$. 
\end{definitionu}

\begin{definitionu} $A$ set $M$ of a topological space $X$ is \emph{closed} if it contains all its accumulation points. It is easy to see that $M$ is closed iff\footnote{\hspace{0.35em}\emph{iff} is the abbreviation for ``if and only if''.} its \emph{complement} $M^{C}=X-M$ is open. Here $A$-$B$ denotes the totality of points $x\in A$ not contained in $B$. If $M\subseteq X$, the intersection of all closed subsets of $X$ which contain $M$ is called the \emph{closure} of $M$ and will be denoted by $M^{a}$ (the superscript ``a'' stands for the first letter of the German: abgeschlossene Hülle). 
\end{definitionu}

Clearly $M^a$ is closed and $M \subseteq M^a$; it is easy to see that $M = M^a$ iff $M$ is closed. 

\yossub{Metric Spaces}

\begin{definitionu} If $X$, $Y$ are sets, we denote by $X \times Y$ the set of all ordered pairs $(x, y)$ where $x \in X$ and $y \in Y$; $X \times Y$ will be called the \emph{Cartesian product} of $X$ and $Y$. $X$ is called a metric space if there is defined a function $d$ with domain $X \times X$ and range in the real number field $\mathbb{R}^{1}$ such that 
$d(x_1,x_2)\geqq0$, and $d(x_1,x_2)=0$ iff $x_1=x_2$; $d(x_1,x_2)=d(x_2,x_1)$; $d(x_1,x_3)\leqq d(x_1,x_2)+d(x_2,x_3)$ (the triangle inequality).
$d$ is called the metric or the distance function of $X$. With each point $x_{0}$ in a metric space $X$ and each positive number $r$, we associate the set $S(x_{0}; r) = \{x \in X; d(x, x_{0}) < r\}$ and call it the open sphere with centre $x_{0}$ and radius $r$. Let us call ``open'' the set $M$ of a metric space $X$ iff, for every point $x_{0} \in M$, $M$ contains a sphere with centre $x_{0}$. Then the totality of such ``open'' sets satisfies the axiom of open sets in the definition of the topological space. 
\end{definitionu}

Hence a metric space $X$ is a topological space. It is easy to see that a point $x_{0}$ of $X$ is an accumulation point of $M$ iff, to every $\varepsilon > 0$, there exists at least one point $m \neq x_{0}$ of $M$ such that $d(m, x_{0}) < \varepsilon$. The $n$-dimensional Euclidean space $\mathbb{R}^{n}$ is a metric space by 
$d(x,y) = \left(\sum\limits_{i=1}^{n}(x_i- y_i)^2\right)^{1/2}$, where $x = (x_1, \ldots, x_n)$ and $y = (y_1, \ldots, y_n)$. 

\yossub{Continuous Mappings}

\begin{definitionu} Let $f: X \to Y$ be a mapping defined on a topological space $X$ into a topological space $Y$. $f$ is called ccontinuous at a point $x_0 \in X$ if to every neighbourhood $U$ of $f(x_0)$ there corresponds a neighbourhood $V$ of $x_0$ such that $f(V) \subseteq U$. The mapping $f$ is said to be continuous if it is continuous at every point of its domain $D(f) = X$. 
\end{definitionu}

\begin{theoremu} Let $X$, $Y$ be topological spaces and $f$ a mapping defined on $X$ into $Y$. Then $f$ is continuous iff the inverse image under $f$ of every open set of $Y$ is an open set of $X$. 
\end{theoremu}

\begin{proof}
 If $f$ is continuous and $U$ an open set of $Y$, then $V = f^{-1}(U)$ is a neighbourhood of every point $x_0 \in X$ such that $f(x_0) \in U$, that is, $V$ is a neighbourhood of every point $x_0$ of $V$. Thus $V$ is an open set of $X$. Let, conversely, for every open set $U \ni f(x_0)$ of $Y$, the set $V = f^{-1}(U)$ be an open set of $X$. Then, by the definition, $f$ is continuous at $x_0 \in X$. 
\end{proof}

\yossub{Compactness}

\begin{definitionu} $A$ system of sets $G_{\alpha}, \alpha \in A$, is called a covering of the set $X$ if $X$ is contained as a subset of the union $\cup_{\alpha \in A} G_{\alpha}$. A subset $M$ of a topological space $X$ is called compact if every system of open sets of $X$ which covers $M$ contains a finite subsystem also covering $M$. 
\end{definitionu}

In view of the preceding theorem, a continuous image of a compact set is also compact. 

\begin{proposition}\label{prop:0.2.Proposition.1} Compact subsets of a topological space are necessarily closed. 
\end{proposition}

\begin{proof}
 Let there be an accumulation point $x_0$ of a compact set $M$ of a topological space $X$ such that $x_0 \in M$. By Hausdorff's axiom of separation, there exist, for any point $m \in M$, disjoint open sets $G_{m,x_0}$ and $G_{x_0,m}$ of $X$ such that $m \in G_{m,x_0}$, $x_0 \in G_{x_0,m}$. The system $\{G_{m,x_0}; m \in M\}$ surely covers $M$. By the compactness of $M$, there exists a finite subsystem $\{G_{m_i,x_0}; i = 1,2,\ldots,n\}$ which covers $M$. Then $\bigcap\limits_{i=1}^{n} G_{x_0,m_i}$ does not intersect $M$. But, since $x_0$ is an accumulation point of $M$, the open set $\bigcap\limits_{i=1}^{n} G_{x_0,m_i} \ni x_0$ must contain a point $m \in M$ distinct from $x_0$. This is a contradiction, and $M$ must be closed. 
\end{proof}

\begin{proposition}\label{prop:0.2.Proposition.2} $A$ closed subset $M_{1}$ of a compact set $M$ of a topological space $X$ is compact. 
\end{proposition}

\begin{proof}
 Let $\{G_{\alpha}\}$ be any system of open sets of $X$ which covers $M_{1}$. $M_{1}$ being closed, $M_{1}^{C}=X-M_{1}$ is an open set of $X$. Since $M_{1}\subseteq M$, the system of open sets $\{G_{\alpha}\}$ plus $M_{1}^{C}$ covers $M$, and since $M$ is compact, a properly chosen finite subsystem $\{G_{\alpha_{i}};i=1,2,\ldots,n\}$ plus $M_{1}^{C}$ surely covers $M$. Thus $\{G_{\alpha_{i}};i=1,2,\ldots,n\}$ covers $M_{1}$. 
\end{proof}

\begin{definitionu} $A$ subset of a topological space is called relatively compact if its closure is compact. A topological space is said to be locally compact if each point of the space has a compact neighbourhood. 
\end{definitionu}

\begin{theoremu} Any locally compact space $X$ can be embedded in another compact space $Y$, having just one more point than $X$, in such a way that the relative topology of $X$ as a subset of $Y$ is just the original topology of $X$. This $Y$ is called a one point compactification of $X$. 
\end{theoremu}

\begin{proof}
 Let $y$ be any element distinct from the points of $X$. Let $\{U\}$ be the class of all open sets in $X$ such that $U^{C} = X- U$ is compact. We remark that $X$ itself $\in \{U\}$. Let $Y$ be the set consisting of the points of $X$ and the point $y$. A set in $Y$ will be called open if either \begin{enumerate}[(i)]
\item it does not contain $y$ and is open as a subset of $X$
\item it does contain $y$ and its intersection with $X$ is a member of $\{U\}$.
\end{enumerate}
 It is easy to see that $Y$ thus obtained is a topological space, and that the relative topology of $X$ coincides with its original topology.

Suppose $\{V\}$ be a family of open sets which covers $Y$. Then there must be some member of $\{V\}$ of the form $U_0 \cup \{y\}$, where $U_0 \in \{U\}$. By the definition of $\{U\}$, $U_0^C$ is compact as a subset of $X$. It is covered by the system of sets $V \cap X$ with $V \in \{V\}$. Thus some finite subsystem: $V_1 \cap X, V_2 \cap X, \ldots, V_n \cap X$ covers $U_0^C$. Consequently, $V_1, V_2, \ldots, V_n$ and $U_0 \cup \{y\}$ cover $Y$, proving that $Y$ is compact. 
\end{proof}

\yossub{Tychonov's Theorem}

\begin{definitionu} Corresponding to each $\alpha$ of an index set $A$, let there be given a topological space $X_{\alpha}$. The Cartesian product $\prod\limits_{\alpha \in A} X_{\alpha}$ is, by definition, the set of all functions $f$ with domain $A$ such that $f(\alpha) \in X_{\alpha}$ for every $\alpha \in A$. We shall write $f = \prod\limits_{\alpha \in A} f(\alpha)$ and call $f(\alpha)$ the $\alpha$-th coordinate of $f$. When $A$ consists of integers $(1, 2, \ldots, n)$, $\prod\limits_{k=1}^{n} X_k$ is usually denoted by $X_1 \times X_2 \times \cdots \times X_n$. We introduce a (weak) topology in the product space $\prod\limits_{\alpha \in A} X_{\alpha}$ by calling ``open'' the sets of the form $\prod\limits_{\alpha \in A} G_{\alpha}$, where the open set $G_{\alpha}$ of $X_{\alpha}$ coincides with $X_{\alpha}$ for all but a finite set of $\alpha$. 
\end{definitionu}

\begin{theoremn}[Tychonov's Theorem] The Cartesian product $X = \prod\limits_{\alpha \in A} X_{\alpha}$ of a system of compact topological spaces $X_{\alpha}$ is also compact. 
\end{theoremn}

\begin{remarku} As is well known, a closed bounded set on the real line $\mathbb{R}^{1}$ is compact with respect to the topology defined by the distance $d(x,y)=|x-y|$ (the Bolzano-Weierstrass theorem). By the way, a subset $M$ of a metric space is said to be bounded, if $M$ is contained in some sphere $S(x_{0},r)$ of the space. Tychonov's theorem implies, in particular, that a parallelopiped: 
\[
- \infty < a _{i} \leqq x _{i} \leqq B_{i} < \infty \quad (i = 1, 2, \dots , n)
\]
\end{remarku}

of the $n$-dimensional Euclidean space $\mathbb{R}^{n}$ is compact. From this we see that $\mathbb{R}^{n}$ is locally compact. 

\begin{proof}[Proof of Tychonov's Theorem.]
A system of sets has the finite intersection property if its every finite subsystem has a non-void intersection. It is easy to see, by taking the complement of the open sets of a covering, that a topological space $X$ is compact iff, for every system $\{M_{\alpha}; \alpha \in A\}$ of its closed subsets with finite intersection property, the intersection $\bigcap\limits_{\alpha \in A} M_{\alpha}^{a}$ is non-void.

Let now a system $\{S\}$ of subsets $S$ of $X = \prod\limits_{\alpha \in A} X_{\alpha}$ have the finite intersection property. Let $\{N\}$ be a system of subsets $N$ of $X$ with the following properties: 

\begin{enumerate}[(i)]
\item $\{S\}$ is a subsystem of $\{N\}$, 

\item $\{N\}$ has the finite intersection property, 

\item $\{N\}$ is maximal in the sense that it is not a proper subsystem of other systems having the finite intersection property and containing $\{S\}$ as its subsystem. 

\end{enumerate}

The existence of such a maximal system $\{N\}$ can be proved by Zorn's lemma or transfinite induction. 

For any set $N$ of $\{N\}$ we define the set $N_{\alpha} = \{f(\alpha); f \in N\} \subseteq X_{\alpha}$. We denote then by $\{N_{\alpha}\}$ the system $\{N_{\alpha}; N \in \{N\}\}$. Like $\{N\}$, $\{N_{\alpha}\}$ enjoys the finite intersection property. Thus, by the compactness of $X_{\alpha}$, there exists at least one point $p_{\alpha} \in X_{\alpha}$ such that $p_{\alpha} \in \bigcap\limits_{N \in \{N\}} N_{\alpha}^{a}$. We have to show that the point $p = \prod\limits_{\alpha \in A} p_{\alpha}$ belongs to the set $\bigcap\limits_{N \in \{N\}} N^{a}$. 

But since $p_{\alpha_0}$ belongs to the intersection $\bigcap\limits_{N \in \{N\}} N_{\alpha_0}^a$, any open set $G_{\alpha_0}$ of $X_{\alpha_0^-}$ which contains $p_{\alpha_0}$ intersects every $N_{\alpha_0} \in \{N_{\alpha_0}\}$. Therefore the open set 
\[
G ^ {(\alpha_{0})} = \{x; x = \prod\limits_{\alpha \in A} x _{\alpha} \text{ with } x _ {\alpha_{0}} \in G _ {\alpha_{0}} \}
\]
of $X$ must intersect every $N$ of $\{N\}$. By the maximality condition (iii) of $\{N\}$, $G^{(\alpha_0)}$ must belong to $\{N\}$. Thus the intersection of a finite number of sets of the form $G^{(\alpha_0)}$ with $\alpha_0 \in A$ must also belong to $\{N\}$ and so such a set intersect every set $N \in \{N\}$. Any open set of $X$ containing $p$ being defined as a set containing such an intersection, we see that $p = \prod\limits_{\alpha \in A} p_\alpha$ must belong to the intersection $\bigcap\limits_{N \in \{N\}} N^a$. 
\end{proof}

\yossub{Urysohn's Theorem}

\begin{propositionu} $A$ compact space $X$ is normal in the sense that, for any disjoint closed sets $F_{1}$ and $F_{2}$ of $X$, there exist disjoint open sets $G_{1}$ and $G_{2}$ such that $F_{1} \subseteq G_{1}, F_{2} \subseteq G_{2}$. 
\end{propositionu}

\begin{proof}
 For any pair $(x, y)$ of points such that $x \in F_1$, $y \in F_2$, there exist disjoint open sets $G(x, y)$ and $G(y, x)$ such that $x \in G(x, y)$, $y \in G(y, x)$. $F_2$ being compact as a closed subset of the compact space $X$, we can, for fixed $x$, cover $F_2$ by a finite number of open sets $G(y_1, x)$, $G(y_2, x), \ldots, G(y_{n(x)}, x)$. Set 
\[
G _{x} = \bigcup\limits_{j=1} ^{n(x)} G (y _{j}, x) \quad \text{ and } \quad G (x) = \bigcap\limits_{j=1} ^{n(x)} G (x, y _{j}). \]

Then the disjoint open sets $G_{x}$ and $G(x)$ are such that $F_{2} \subseteq G_{x}, x \in G(x) \cdot F_{1}$ being compact as a closed subset of the compact space $X$, we can cover $F_{1}$ by a finite number of open sets $G(x_{1}), G(x_{2}), \ldots, G(x_{k})$. Then 
\[
G _{1} = \bigcup\limits_{j=1} ^{k} G (x _{j}) \quad \text{ and } \quad G _{2} = \bigcap\limits_{j=1} ^{k} G _ {x _{j}}
\]
satisfy the condition of the proposition. 
\end{proof}

\begin{corollaryu} $A$ compact space $X$ is regular in the sense that, for any non-void open set $G_1'$ of $X$, there exists a non-void open set $G_2'$ such that $(G_2')^a \subseteq G_1'$. 
\end{corollaryu}

\begin{proof}
 Take $F_{1} = (G_{1}')^{C}$ and $F_{2} = \{x\}$ where $x \in G_{1}'$. We can then take for $G_{2}'$ the open set $G_{2}$ obtained in the preceding proposition. 
\end{proof}

\begin{theoremn}[Urysohn's Theorem] Let $A, B$ be disjoint closed sets in a normal space $X$. Then there exists a real-valued continuous function $f(t)$ on $X$ such that 
\[
0 \leqq f (t) \leqq 1 \text{ on } X, \text{ and } f (t) = 0 \text{ on } A, f (t) = 1 \text{ on } B.
\]
\end{theoremn}

\begin{proof}
 We assign to each rational number $r = k / 2^n$ ( $k = 0, 1, \ldots, 2^n$), an open set $G(r)$ such that \begin{enumerate}[(i)]
\item $A \subseteq G(0)$, $B = G(1)^c$
\item $G(r)^a \subseteq G(r')$ whenever $r < r'$.
\end{enumerate}
 The proof is obtained by induction with respect to $n$.

For $n = 0$, there exist, by the normality of the space $X$, disjoint open sets $G_0$ and $G_1$ with $A \subseteq G_0$, $B \subseteq G_1$. We have only to set $G_0 = G(0)$. Suppose that $G(r)'s$ have been constructed for $r$ of the form $k/2^{n-1}$ in such a way that condition (ii) is satisfied. Next let $k$ be an odd integer $> 0$. Then, since $(k-1)/2^n$ and $(k+1)/2^n$ are of the form $k'/2^{n-1}$ with $0 \leqq k' \leqq 2^{n-1}$, we have $G((k-1)/2^n)^a \subseteq G((k+1)/2^n)$. Hence, by the normality of the space $X$, there exists an open set $G$ which satisfies $G((k-1)/2^n)^a \subseteq G$, $G^a \subseteq G((k+1)/2^n)$. If we set $G(k/2^n) = G$, the induction is completed. 

Define $f(t)$ by 
\[
f (t) = 0 \text{ on } G (0), \text{ and } f (t) = \sup\limits _{t \in G (r)} r \text{ whenever } t \in G (0) ^{C}.
\]

Then, by (i), $f(t) = 0$ on $A$ and $f(t) = 1$ on $B$. We have to prove the continuity of $f$. For any $t_0 \in X$ and positive integer $n$, we take $r$ with $f(t_0) < r < f(t_0) + 2^{-n-1}$. Set $G = G(r) \cap G(r- 2^{-n})^{aC}$ (we set, for convention, $G(s) = \emptyset$ if $s < 0$ and $G(s) = X$ if $s > 1$). The open set $G$ contains $t_0$. For, $f(t_0) < r$ implies $t_0 \in G(r)$, and $(r- 2^{-n-1}) < f(t_0)$ implies $t_0 \in G(r- 2^{-n-1})^C \subseteq G(r- 2^{-n})^{aC}$. Now $t \in G$ implies $t \in G(r)$ and so $f(t) \leqq r$; similarly $t \in G$ implies $t \in G(r- 2^{-n})^{aC} \subseteq G(r- 2^{-n})^C$ so that $r- 2^{-n} \leqq f(t)$. Therefore we have proved that $|f(t)- f(t_0)| \leqq 1/2^n$ whenever $t \in G$. 
\end{proof}
\yossub{The Stone-Weierstrass Theorem}

\begin{theoremn}[Weierstrass' Polynomial Approximation Theorem] Let $f(x)$ be a real-valued (or complex-valued) continuous function on the closed interval [$0$, $1$]. Then there exists a sequence of polynomials $P_{n}(x)$ which converges, as $n \to \infty$, to $f(x)$ uniformly on [$0$, $1$]. According to S. \textsc{Bernstein}, we may take 
\begin{equation}\label{eq:0.2.1}
P _{n} (x) = \sum\limits_{p=0} ^{n} n C _{p} f (p / n) x ^{p} (1- x) ^{n-p}.
\end{equation}
\end{theoremn}

\begin{proof}
 Differentiating $(x + y)^n = \sum\limits_{p=0}^{n} {}_n C_p x^p y^{n-p}$ with respect to $x$ and multiplying by $x$, we obtain $nx(x + y)^{n-1} = \sum\limits_{p=0}^{n} p {}_n C_p x^p y^{n-p}$. Similarly, by differentiating the first expression twice with respect to $x$ and multiplying by $x^2$, we obtain $n(n-1)x^2(x + y)^{n-2} = \sum\limits_{p=0}^{n} p(p-1) {}_n C_p x^p y^{n-p}$. Thus, if we set 
\begin{equation}\label{eq:0.2.2}
r _{p} (x) = _{n} C _{p} x ^{p} (1- x) ^{n-p},
\end{equation}
we have 
\begin{equation}\label{eq:0.2.3}
\sum\limits_{p=0} ^{n} r _{p} (x) = 1, \sum\limits_{p=0} ^{n} p r _{p} (x) = n x, \sum\limits_{p=0} ^{n} p (p- 1) r _{p} (x) = n (n- 1) x ^{2}.
\end{equation}

Hence 
\begin{equation}\label{eq:0.2.4}
\begin{aligned} \sum\limits_{p=0} ^{n} (p- n x) ^{2} r _{p} (x) & = n ^{2} x ^{2} \sum\limits_{p=0} ^{n} r _{p} (x)- 2 n x \sum\limits_{p=0} ^{n} p r _{p} (x) + \sum\limits_{p=0} ^{n} p ^{2} r _{p} (x) \\ & = n ^{2} x ^{2}- 2 n x \cdot n x + (n x + n (n- 1) x ^{2}) \\ & = n x (1- x). \end{aligned}
\end{equation}

We may assume that $|f(x)| \leqq M < \infty$ on [$0$, $1$]. By the uniform continuity of $f(x)$, there exists, for any $\varepsilon > 0$, a $\delta > 0$ such that 
\begin{equation}\label{eq:0.2.5}
\left| f (x)- f \left(x ^{\prime}\right) \right| < \varepsilon \text{ whenever } \left| x- x ^{\prime} \right| < \delta .
\end{equation}

We have, by \eqref{eq:0.2.3}, 
\[
\begin{aligned}
&\left| f (x)- \sum\limits_{p=0} ^{n} f (p / n) r _{p} (x) \right| = \left| \sum\limits_{p=0} ^{n} (f (x)- f (p / n)) r _{p} (x) \right| \leqq \left| \sum\limits_{| p- n x | \leqq \delta n} \right| + \left| \sum\limits_{| p- n x | > \delta n} \right|.
\end{aligned}
\]

For the first term on the right, we have, by $r_{p}(x) \geqq 0$ and \eqref{eq:0.2.3}, 
\[
\left| \sum\limits_{| p- n x | \leqq \delta n} \right| \leqq \varepsilon \sum\limits_{p=0} ^{n} r _{p} (x) = \varepsilon .
\]

For the second term on the right, we have, by \eqref{eq:0.2.4} and $|f(x)| \leqq M$, 
\[ \begin{aligned}\left| \sum\limits_{| p- n x | > \delta n} \right|&\leqq 2 M \sum\limits_{| p- n x | > \delta n} r _{p} (x) \leqq \frac {2 M}{n ^{2} \delta^{2}} \sum\limits_{p=0} ^{n} (p- n x) ^{2} r _{p} (x)\\&= \frac {2 M x (1- x)}{n \delta^{2}} \leqq \frac {M}{2 \delta^{2} n} \rightarrow 0 (\text{ as } n \rightarrow \infty).\end{aligned}
\]
\end{proof}
\begin{theoremn}[The Stone-Weierstrass Theorem] Let $X$ be a compact space and $C(X)$ the totality of real-valued continuous functions on $X$. Let a subset $B$ of $C(X)$ satisfy the three conditions: \begin{enumerate}[(i)]
\item if $f, g \in B$, then the function product $f \cdot g$ and linear combinations $\alpha f + \beta g$, with real coefficients $\alpha, \beta$, belong to $B$,
\item the constant function $1$ belongs to $B$
\item the uniform limit $f_{\infty}$ of any sequence $\{f_n\}$ of functions $\in B$ also belongs to $B$.
\end{enumerate}
 Then $B = C(X)$ iff $B$ separates the points of $X$, i.e. iff, for every pair $(s_1, s_2)$ of distinct points of $X$, there exists a function $x$ in $B$ which satisfies $x(s_1) \neq x(s_2)$. 
\end{theoremn}

\begin{proof}
 The necessity is clear, since a compact space is normal and so, by Urysohn's theorem, there exists a real-valued continuous function $x$ such that $x(s_1) \neq x(s_2)$. 

To prove the sufficiency, we introduce the lattice notations: 
\[
(f \vee g) (s) = \max (f (s), g (s)), (f \wedge g) (s) = \min (f (s), g (s)), | f | (s) = | f (s) |.
\]

By the preceding theorem, there is a sequence $\{P_n\}$ of polynomials such that 
\[
\left| \left| t \right|- P _{n} (t) \right| < 1 / n \quad \text{ for } \quad- n \leqq t \leqq n.
\]

Hence $\left|\left|f(s)\right|- P_n(f(s))\right| < 1/n$ if $-n \leqq f(s) \leqq n$. This proves, by (iii), that $|f| \in B$ if $f \in B$, because any function $f(s) \in B \subseteq C(X)$ is bounded on the compact space $X$. Thus, by 
\[
f \vee g = \frac {f + g}{2} + \frac {| f- g |}{2} \quad \text{ and } \quad f \wedge g = \frac {f + g}{2}- \frac {| f- g |}{2},
\]
we see that $B$ is closed under the lattice operations $\vee$ and $\wedge$. 

Let $h \in C(X)$ and $s_1, s_2 \in X$ be arbitrarily given such that $s_1 \neq s_2$. Then we can find an $f_{s_1s_2} \in B$ with $f_{s_1s_2}(s_1) = h(s_1)$ and $f_{s_1s_2}(s_2) = h(s_2)$. To see this, let $g \in B$ be such that $g(s_1) \neq g(s_2)$, and take real numbers $\alpha$ and $\beta$ so that $f_{s_1s_2} = \alpha g + \beta$ satisfies the conditions: $f_{s_1s_2}(s_1) = h(s_1)$ and $f_{s_1s_2}(s_2) = h(s_2)$. 

Given $\varepsilon > 0$ and a point $t \in X$. Then, for each $s \in X$, there is a neighbourhood $U(s)$ of $s$ such that $f_{st}(u) > h(u)- \varepsilon$ whenever $u \in U(s)$. Let $U(s_1), U(s_2), \ldots, U(s_n)$ cover the compact space $X$ and define 
\[
f _{t} = f _ {s _{1} t} \vee \dots \vee f _ {s _{n} t}.
\]

Then $f_{t} \in B$ and $f_{t}(u) > h(u)- \varepsilon$ for all $u \in X$. We have, by $f_{s,t}(t) = h(t)$, $f_{t}(t) = h(t)$. Hence there is a neighbourhood $V(t)$ of $t$ such that $f_{t}(u) < h(u) + \varepsilon$ whenever $u \in V(t)$. Let $V(t_{1}), V(t_{2}), \ldots, V(t_{k})$ cover the compact space $X$, and define 
\[
f = f _ {t _{1}} \wedge \dots \wedge f _ {t _{k}}.
\]

Then $f \in B$ and $f(u) > h(u)- \varepsilon$ for all $u \in X$, because $f_{t_{j}}(u) > h(u)- \varepsilon$ for $u \in X$. Moreover, we have, for an arbitrary point $u \in X$, say $u \in V(t_i)$, $f(u) \leqq f_{t_{i}}(u) < h(u) + \varepsilon$. 

Therefore we have proved that $|f(u)- h(u)| < \varepsilon$ on $X$. 

We have incidentally proved the following two corollaries. 
\end{proof}

\begin{corollary}[\textsc{Kakutani}-\textsc{Krein}]\label{cor:0.2.Corollary.1} Let $X$ be a compact space and $C(X)$ the totality of real-valued continuous functions on $X$. Let a subset $B$ of $C(X)$ satisfy the conditions: \begin{enumerate}[(i)]
\item if $f, g \in B$, then $f \vee g, f \wedge g$ and the linear combinations $\alpha f + \beta g$, with real coefficients $\alpha, \beta$, belong to $B$,
\item the constant function $1$ belongs to $B$
\item the uniform limit $f_{\infty}$ of any sequence $\{f_n\}$ of functions $\in B$ also belongs to $B$.
\end{enumerate}
 Then $B = C(X)$ iff $B$ separates the points of $X$. 
\end{corollary}

\begin{corollary}\label{cor:0.2.Corollary.2} Let $X$ be a compact space and $C(X)$ be the totality of complex-valued continuous functions on $X$. Let a subset $B$ of $C(X)$ satisfy the conditions: \begin{enumerate}[(i)]
\item if $f, g \in B$, then the function product $f \cdot g$ and the linear combinations $\alpha f + \beta g$, with complex coefficients $\alpha, \beta$, belong to $B$,
\item the constant function $1$ belongs to $B$
\item the uniform limit $f_{\infty}$ of any sequence $\{f_n\}$ of functions $\in B$ also belongs to $B$.
\end{enumerate}
 Then $B = C(X)$ iff $B$ satisfies the conditions: \begin{enumerate}[(i)]
\setcounter{enumi}{3}
\item $B$ separates points of $X$
\item if $f(s) \in B$, then its complex conjugate function $\bar{f}(s)$ also belongs to $B$.
\end{enumerate}
\end{corollary}

\begin{theoremn}[Weierstrass' Trigonometric Approximation Theorem] Let $X$ be the circumference of the unit circle of $\mathbb{R}^{2}$. It is a compact space by the usual topology, and a complex-valued continuous function on $X$ is represented by a continuous function $f(x),-\infty < x < \infty$, of period $2\pi$. If we take, in the above Corollary~\ref{cor:0.2.Corollary.2}, for $B$ the set of all functions representable by linear combinations, with complex coefficients of the trigonometric functions 
\[
e ^{inx} (n = 0, \pm 1, \pm 2, \dots)
\]

and by those functions obtainable as the uniform limit of such linear combinations, we obtain Weierstrass' trigonometric approximation theorem: Any complex-valued continuous function $f(x)$ with period $2\pi$ can be approximated uniformly by a sequence of trigonometric polynomials of the form $\sum\limits_{n} c_n e^{inx}$. 
\end{theoremn}

\yossub{Completeness}

A sequence $\{x_{n}\}$ of elements in a metric space $X$ converges to a limit point $x\in X$ iff $\lim\limits_{n\to \infty}d(x_n,x) = 0$. By the triangle inequality $\tilde{d}(x_n,x_m)\leqq d(x_n,x) + d(x,x_m)$, we see that a convergent sequence $\{x_{n}\}$ in $X$ satisfies Cauchy's convergence condition 
\begin{equation}\label{eq:0.2.1a}
\lim\limits _{n, m \rightarrow \infty} d \left(x _{n}, x _{m}\right) = 0.
\end{equation}

\begin{definitionu} Any sequence $\{x_{n}\}$ in a metric space $X$ satisfying the above condition \eqref{eq:0.2.1a} is called a Cauchy sequence. A metric space $X$ is said to be complete if every Cauchy sequence in it converges to a limit point $\in X$. 
\end{definitionu}

It is easy to see, by the triangle inequality, that the limit point of $\{x_{n}\}$, if it exists, is uniquely determined. 

\begin{definitionu} $A$ subset $M$ of a topological space $X$ is said to be non-dense in $X$ if the closure $M^{a}$ does not contain a non-void open set of $X$. $M$ is called dense in $X$ if $M^{a} = X$. $M$ is said to be of the first category if $M$ is expressible as the union of a countable number of sets each of which is non-dense in $X$; otherwise $M$ is said to be of the second category. 
\end{definitionu}

\yossub{Baire's Category Argument}

\begin{theoremn}[The Baire-Hausdorff Theorem] $A$ non-void complete metric space is of the second category. 
\end{theoremn}

\begin{proof}
 Let $\{M_{n}\}$ be a sequence of closed sets whose union is a complete metric space $X$. Assuming that no $M_{n}$ contains a non-void open set, we shall derive a contradiction. Thus $M_{1}^{C}$ is open and $M_{1}^{Ca}=X$, hence $M_{1}^{C}$ contains a closed sphere $S_{1}=\{x;d(x_{1},x)\leqq r_{1}\}$ whose centre $x_{1}$ may be taken arbitrarily near to any point of $X$. We may assume that $0<r_{1}<1/2$. By the same argument, the open set $M_{2}^{C}$ contains a closed sphere 

$S_{2} = \{x; d(x_{2}, x) \leqq r_{2}\}$ contained in $S_{1}$ and such that $0 < r_{2} < 1/2^{2}$. By repeating the same argument, we obtain a sequence $\{S_{n}\}$ of closed spheres $S_{n} = \{x; d(x_{n}, x) \leqq r_{n}\}$ with the properties: 
\[
0 < r _{n} < 1 / 2 ^{n}, S _{n+1} \subseteq S _{n}, S _{n} \cap M _{n} = \emptyset \quad (n = 1, 2, \dots). \]

The sequence $\{x_{n}\}$ of the centres forms a Cauchy sequence, since, for any $n < m$, $x_{m} \in S_{n}$ so that $d(x_{n}, x_{m}) \leqq r_{n} < 1/2^{n}$. Let $x_{\infty} \in X$ be the limit point of this sequence $\{x_{n}\}$. The completeness of $X$ guarantees the existence of such a limit point $x_{\infty}$. By $d(x_{n}, x_{\infty}) \leqq d(x_{n}, x_{m}) + d(x_{m}, x_{\infty}) \leqq r_{n} + d(x_{m}, x_{\infty}) \to r_{n}$ (as $m \to \infty$), we see that $x_{\infty} \in S_{n}$ for every $n$. Hence $x_{\infty}$ is in none of the sets $M_{n}$, and hence $x_{\infty}$ is not in the union $\bigcup\limits_{n=1}^{\infty} M_{n} = X$, contrary to $x_{\infty} \in X$. 
\end{proof}

\begin{theoremn}[Baire's Theorem 1] Let $M$ be a set of the first category in a compact topological space $X$. Then the complement $M^C = X- M$ is dense in $X$. 
\end{theoremn}

\begin{proof}
 We have to show that, for any non-void open set $G$, $M^C$ intersects $G$. Let $M = \bigcup\limits_{n=1}^{\infty} M_n$ where each $M_n$ is a non-dense closed set. Since $M_1 = M_1^a$ is non-dense, the open set $M_1^C$ intersects $G$. Since $X$ is regular as a compact space, there exists a non-void open set $G_1$ such that $G_1^a \subseteq G \cap M_1^C$. Similarly, we can choose a non-void open set $G_2$ such that $G_2^a \subseteq G_1 \cap M_2^C$. Repeating the process, we obtain a sequence of non-void open sets $\{G_n\}$ such that 
\[
G _{n+1} ^{a} \subseteq G _{n} \cap M _{n+1} ^{C} \quad (n = 1, 2, \dots). \]
The sequence of closed sets $\{G_n^a\}$ enjoys, by the monotony in $n$, the finite intersection property. Since $X$ is compact, there is an $x \in X$ such that $x \in \bigcap\limits_{n=1}^{\infty} G_n^a$. $x \in G_1^a$ implies $x \in G$, and from $x \in G_{n+1}^a \subseteq G_n \cap M_{n+1}^C$ ( $n = 0, 1, 2, \ldots; G_0 = G$), we obtain $x \in \bigcap\limits_{n=1}^{\infty} M_n^c = M^c$. Therefore we have proved that $G \cap M^c$ is non-void. 
\end{proof}

\begin{theoremn}[Baire's Theorem 2] Let $\{x_{n}(t)\}$ be a sequence of real-valued continuous functions defined on a topological space $X$. Suppose that a finite limit: 
\[
\lim\limits _{n \rightarrow \infty} x _{n} (t) = x (t)
\]

exists at every point $t$ of $X$. Then the set of points at which the function $x$ is discontinuous constitutes a set of the first category. 
\end{theoremn}

\begin{proof}
 We denote, for any set $M$ of $X$, by $M^{i}$ the union of all the open sets contained in $M$; $M^{i}$ will be called the interior of $M$. 

 Put $P_{m} (\varepsilon) = \{t \in X; | x (t)- x _{m} (t) | \leqq \varepsilon , \varepsilon > 0 \}, G (\varepsilon) = \bigcup\limits_{m=1} ^{\infty} P _{m} ^{i} (\varepsilon). $ Then we can prove that $C = \bigcap\limits_{n=1}^{\infty} G(1/n)$ coincides with the set of all points at which $x(t)$ is continuous. Suppose $x(t)$ is continuous at $t = t_0$. 

We shall show that $t_0 \in \bigcap\limits_{n=1}^{\infty} G(1/n)$. Since $\lim\limits_{n \to \infty} x_n(t) = x(t)$, there exists an $m$ such that $|x(t_0)- x_m(t_0)| \leqq \varepsilon/3$. By the continuity of $x(t)$ and $x_m(t)$ at $t = t_0$, there exists an open set $U_{t_0} \ni t_0$ such that $|x(t)- x(t_0)| \leqq \varepsilon/3$ $|x_m(t)- x_m(t_0)| \leqq \varepsilon/3$ whenever $t \in U_{t_0}$. Thus $t \in U_{t_0}$ implies $|x(t)- x_m(t)| \leqq |x(t)- x(t_0)| + |x(t_0)- x_m(t_0)| + |x_m(t_0)- x_m(t)| < \varepsilon$, which proves that $t_0 \in P_m^i(\varepsilon)$ and so $t_0 \in G(\varepsilon)$. Since $\varepsilon > 0$ was arbitrary, we must have $t_0 \in \bigcap\limits_{n=1}^{\infty} G(1/n)$. 

Let, conversely, $t_0 \in \bigcap\limits_{n=1}^{\infty} G(1/n)$. Then, for any $\varepsilon > 0$, $t_0 \in G(\varepsilon/3)$ and so there exists an $m$ such that $t_0 \in P_m^i(\varepsilon/3)$. Thus there is an open set $U_{t_0} \ni t_0$ such that $t \in U_{t_0}$ implies $|x(t)- x_m(t)| \leqq \varepsilon/3$. Hence, by the continuity of $x_m(t)$ and the arbitrariness of $\varepsilon > 0$, $x(t)$ must be continuous at $t = t_0$. 

After these preparations, we put 
\[
F _{m} (\varepsilon) = \{t \in X; | x _{m} (t)- x _{m+k} (t) | \leqq \varepsilon \quad (k = 1, 2, \dots) \}.
\]
This is a closed set by the continuity of the $x_{n}(t)$'s. We have $X = \bigcup\limits_{m=1}^{\infty} F_{m}(\varepsilon)$ by $\lim\limits_{n \to \infty} x_{n}(t) = x(t)$. Again by $\lim\limits_{n \to \infty} x_{n}(t) = x(t)$, we have $F_{m}(\varepsilon) \subseteq P_{m}(\varepsilon)$. Thus $F_{m}^{i}(\varepsilon) \subseteq P_{m}^{i}(\varepsilon)$ and so $\bigcup\limits_{m=1}^{\infty} F_{m}^{i}(\varepsilon) \subseteq G(\varepsilon)$. On the other hand, for any closed set $F$, $(F- F^{i})$ is a non-dense closed set. Thus $X- \bigcup\limits_{m=1}^{\infty} F_{m}^{i}(\varepsilon) = \bigcup\limits_{m=1}^{\infty} (F_{m}(\varepsilon)- F_{m}^{i}(\varepsilon))$ is a set of the first category. Thus its subset $G(\varepsilon)^{C} = X- G(\varepsilon)$ is also a set of the first category. Therefore the set of all the points of discontinuity of the function $x(t)$, which is expressible as $X- \bigcap\limits_{n=1}^{\infty} G(1/n) = \bigcup\limits_{n=1}^{\infty} G(1/n)^{C}$, is a set of the first category. 
\end{proof}

\begin{theoremu} $A$ subset $M$ of a complete metric space $X$ is relatively compact iff it is totally bounded in the sense that, for every $\varepsilon > 0$, there exists a finite system of points $m_{1}, m_{2}, \ldots, m_{n}$ of $M$ such that every point $m$ of $M$ has a distance $< \varepsilon$ from at least one of $m_{1}, m_{2}, \ldots, m_{n}$. In other words, $M$ is totally bounded if, for every $\varepsilon > 0$, $M$ can be covered by a finite system of spheres of radii $< \varepsilon$ and centres $\in M$. 
\end{theoremu}

\begin{proof}
 Suppose $M$ is not totally bounded. Then there exist a positive number $\varepsilon$ and an infinite sequence $\{m_{n}\}$ of points $\in M$ such that $d(m_{i}, m_{j}) \geqq \varepsilon$ for $i \neq j$. Then, if we cover the compact set $M^{a}$ by a system of open spheres of radii $< \varepsilon$, no finite subsystem of this system can cover $M^{a}$. For, this subsystem cannot cover the infinite subset $\{m_{i}\} \subseteq M \subseteq M^{a}$. Thus a relatively compact subset of $X$ must be totally bounded. 

Suppose, conversely, that $M$ is totally bounded as a subset of a complete metric space $X$. Then the closure $M^a$ is complete and is totally bounded with $M$. We have to show that $M^a$ is compact. To this purpose, we shall first show that any infinite sequence $\{p_n\}$ of $M^a$ contains a subsequence $\{p_{(n')}\}$ which converges to a point of $M^a$. Because of the total boundedness of $M$, there exist, for any $\varepsilon > 0$, a point $q_e \in M^a$ and a subsequence $\{p_{n'}\}$ of $\{p_n\}$ such that $d(p_{n'}, q_e) < \varepsilon / 2$ for $n = 1, 2, \ldots$; consequently, $d(p_{n'}, p_{m'}) \leqq d(p_{n'}, q_e) + d(q_e, p_{m'}) < \varepsilon$ for $n, m = 1, 2, \ldots$. We set $\varepsilon = 1$ and obtain the sequence $\{p_{i'}\}$, and then apply the same reasoning as above with $\varepsilon = 2^{-1}$ to this sequence $\{p_{i'}\}$. We thus obtain a subsequence $\{p_{n''}\}$ of $\{p_{n'}\}$ such that 
\[
d \left(p _ {n ^{\prime}}, p _ {m ^{\prime}}\right) < 1, d \left(p _ {n ^{\prime \prime}}, p _ {m ^{\prime \prime}}\right) < 1 / 2 \quad (n, m = 1, 2, \dots).
\]

By repeating the process, we obtain a subsequence $\{\phi_{n^{(k+1)}}\}$ of the sequence $\{\phi_{n^{(k)}}\}$ such that 
\[
d \left(\phi_ {n ^{(k+1)}}, \phi_ {m ^{(k+1)}}\right) < 1 / 2 ^{k} \quad (n, m = 1, 2, \dots).
\]

Then the subsequence $\{\phi_{(n)'}\}$ of the original sequence $\{\phi_n\}$, defined by the diagonal method: 
\[
\phi_ {(n) ^{\prime}} = \phi_ {n ^{(n)}},
\]
surely satisfies $\lim\limits_{n,m\to \infty}d(p_{(n)'}p_{(m)'}) = 0$. Hence, by the completeness of $M^a$, there must exist a point $p\in M^a$ such that $\lim\limits_{n\to \infty}d(p_{(n)'}p) = 0$. 

We next show that the set $M^{a}$ is compact. We remark that there exists a countable family $\{F\}$ of open sets $F$ of $X$ such that, if $U$ is any open set of $X$ and $x \in U \cap M^{a}$, there is a set $F \in \{F\}$ for which $x \in F \subseteq U$. This may be proved as follows. $M^{a}$ being totally bounded, $M^{a}$ can be covered, for any $\varepsilon > 0$, by a finite system of open spheres of radii $\varepsilon$ and centres $\in M^{a}$. Letting $\varepsilon = 1, 1/2, 1/3, \ldots$ and collecting the countable family of the corresponding finite systems of open spheres, we obtain the desired family $\{F\}$ of open sets. 

Let now $\{U\}$ be any open covering of $M^a$. Let $\{F^*\}$ be the subfamily of the family $\{F\}$ defined as follows: $F \subseteq \{F^*\}$ iff $F \subseteq \{F\}$ and there is some $U \in \{U\}$ with $F \subseteq U$. By the property of $\{F\}$ and the fact that $\{U\}$ covers $M^a$, we see that this countable family $\{F^*\}$ of open sets covers $M^a$. Now let $\{U^*\}$ be a subfamily of $\{U\}$ obtained by selecting just one $U \in \{U\}$ such that $F \subseteq U$, for each $F \in \{F^*\}$. Then $\{U^*\}$ is a countable family of open sets which covers $M^a$. We have to show that some finite subfamily of $\{U^*\}$ covers $M^a$. Let the sets in $\{U^*\}$ be indexed as $U_1, U_2, \ldots$. Suppose that, for each $n$, the finite union $\bigcup\limits_{j=1}^{n} U_j$ fails to cover $M^a$. Then there is some point $x_n \in \left(M- \bigcup\limits_{k=1}^{n} U_k\right)$. By what was proved above, the sequence $\{x_{n}\}$ contains a subsequence $\{x_{(n)}\}$ which converges to a point, say $x_{\infty}$, in $M^{a}$. Then $x_{\infty} \in U_{N}$ for some index $N$, and so $x_{n} \in U_{N}$ for infinitely many values of $n$, in particular for an $n > N$. This contradicts the fact that $x_{n}$ was chosen so that $x_{n} \in \left(M- \bigcup\limits_{k=1}^{n} U_{k}\right)$. Hence we have proved that $M^{a}$ is compact. 
\end{proof}


\yossection{3}{Measure Spaces}
\label{sec:0.3}

\yossub{Measures}

\begin{definitionu} Let $S$ be a set. A pair $(S, \mathfrak{B})$ is called a \emph{$\sigma$-ring} or a \emph{$\sigma$-additive family} of sets $\subseteq S$ if $\mathfrak{B}$ is a family of subsets of $S$ such that 
\begin{align}
S& \in \mathfrak{B}, \label{eq:0.3.1} \\
B& \in \mathfrak{B} \text{ implies } B ^{C} = (S- B) \in \mathfrak{B}, \label{eq:0.3.2} \\
B_{j}& \in \mathfrak{B} (j = 1, 2, \dots) \text{ implies that } \bigcup\limits_{j=1} ^{\infty} B_{j} \in \mathfrak{B} (\sigma \text{-additivity }). \label{eq:0.3.3}
\end{align}
\end{definitionu}

Let $(S, \mathfrak{B})$ be a $\sigma$-ring of sets $\subseteq S$. Then a triple $(S, \mathfrak{B}, m)$ is called a \emph{measure space} if $m$ is a non-negative, $\sigma$-additive measure defined on $\mathfrak{B}$: 
\begin{equation}
 m(B) \geqq 0 \text{ for every } B \in \mathfrak{B}, \label{eq:0.3.4} \\
\end{equation}
\begin{equation}\label{eq:0.3.5}
\parbox[b]{0.92\linewidth}{\centering
$m\left(\sum\limits_{j=1}^{\infty}B_j\right)=\sum\limits_{j=1}^{\infty}m(B_j)$ for any disjoint sequence $\{B_j\}$ of sets $\in\mathfrak{B}$\par
(\emph{countable- or $\sigma$-additivity of $m$}),}
\end{equation}
\par\medskip\refstepcounter{equation}\label{eq:0.3.6}
\noindent
\parbox[b]{\linewidth}{\setlength{\parindent}{1.2em}\setlength{\parfillskip}{0pt}\indent
$S$ is expressible as a countable union of sets $B_j\in\mathfrak{B}$ such that $m(B_j)<\infty$ $(j=1,2,\ldots)$ (\emph{$\sigma$-finiteness} of the measure space $(S,\mathfrak{B},m)$).\hfill(\theequation)}
\par\medskip
This value $m(B)$ is called the \emph{$m$-measure} of the set $B$. 

\yossub{Measurable Functions}

\begin{definitionu} $A$ real- (or complex-) valued function $x(s)$ defined on $S$ is said to be \emph{$\mathfrak{B}$-measurable} or, in short, \emph{measurable} if the following condition is satisfied: 

\par\medskip\refstepcounter{equation}\label{eq:0.3.7}
\noindent
\parbox[b]{\linewidth}{\setlength{\parindent}{1.2em}\setlength{\parfillskip}{0pt}\indent
For any open set $G$ of the real line $\mathbb{R}^{1}$ (or complex plane $\mathbb{C}^{1}$), the set $\{s; x(s)\in G\}$ belongs to $\mathfrak{B}$.\hfill(\theequation)}
\par\medskip
It is permitted that $x(s)$ takes the value $\infty$.
\end{definitionu}

\begin{definitionu} $A$ property $P$ pertaining to points $s$ of $S$ is said to hold \emph{$m$-almost everywhere} or, in short \emph{$m$-a.e.}, if it holds except for those $s$ which form a set $\in \mathfrak{B}$ of $m$-measure zero. 
\end{definitionu}

A real- (or complex-) valued function $x(s)$ defined $m$-a.e. on $S$ and satisfying condition \eqref{eq:0.3.7} shall be called a $\mathfrak{B}$-measurable function defined $m$-a.e. on $S$ or, in short, a $\mathfrak{B}$-measurable function. 

\begin{theoremn}[Egorov's Theorem] If $B$ is a $\mathfrak{B}$-measurable set with $m(B) < \infty$ and if $\{f_n(s)\}$ is a sequence of $\mathfrak{B}$-measurable functions, finite $m$-a.e. on $B$, that converges $m$-a.e. on $B$ to a finite $\mathfrak{B}$-measurable function $f(s)$, then there exists, for each $\varepsilon > 0$, a subset $E$ of $B$ such that $m(B- E) \leqq \varepsilon$ and on $E$ the convergence of $f_n(s)$ to $f(s)$ is uniform. 
\end{theoremn}

\begin{proof}
 By removing from $B$, if necessary, a set of $m$-measure zero, we may suppose that on $B$, the functions $f_{n}(s)$ are everywhere finite, and converge to $f(s)$ on $B$. 


The set $B_{n} = \bigcap\limits_{k=n+1}^{\infty}\{s\in B; |f(s)- f_{k}(s)| < \varepsilon\}$ is $\mathfrak{B}$-measurable and $B_{n} \subseteq B_{k}$ if $n < k$. Since $\lim\limits_{n\to\infty}f_{n}(s) = f(s)$ on $B$, we have $B = \bigcup\limits_{n=1}^{\infty}B_{n}$. Thus, by the $\sigma$-additivity of the measure $m$, we have 
\[
\begin{aligned}
&m(B)  = m \{B_{1} + (B_{2}- B_{1}) + (B_{3}- B_{2}) + \dots \} = m(B_{1}) + m(B_{2}- B_{1}) + m(B_{3}- B_{2}) + \dots \\
&= m(B_{1}) + (m(B_{2})- m(B_{1})) + (m(B_{3})- m(B_{2})) + \dots = \lim\limits _{n \to \infty} m(B_{n}).
\end{aligned}
\]

Hence $\lim\limits_{n\to \infty}m(B- B_n) = 0$, and therefore, from a sufficiently large $k_{0}$ on, $m(B- B_k) < \eta$ where $\eta$ is any given positive number. 

Thus there exist, for any positive integer $k$, a set $C_k \subseteq B$ such that $m(C_k) \leqq \varepsilon / 2^k$ and an index $N_k$ such that 
\[
\left| f (s)- f _{n} (s) \right| < 1 / 2 ^{k} \text{ for } n > N _{k} \text{ and for } s \in B- C _{k}.
\]

Let us set $E = B- \bigcup\limits_{k=1}^{\infty} C_k$. Then we find 
\[
m(B- E) \leqq \sum\limits_{k=1} ^{\infty} m(C _{k}) \leqq \sum\limits_{k=1} ^{\infty} \varepsilon / 2 ^{k} = \varepsilon ,
\]
and the sequence $f_{n}(s)$ converges uniformly on $E$.
\end{proof}

\yossub{Integrals}

\begin{definitionu} $A$ real- (or complex-) valued function $x(s)$ defined on $S$ is said to be finitely-valued if it is a finite non-zero constant on each of a finite number, say $n$, of disjoint $\mathfrak{B}$-measurable sets $B_j$ and equal to zero on $S- \bigcup\limits_{j=1}^{n} B_j$. Let the value of $x(s)$ on $B_j$ be denoted by $x_j$. Then $x(s)$ is $m$-integrable or, in short, integrable over $S$ if $\sum\limits_{j=1}^{n} |x_j| m(B_j) < \infty$, and the value $\sum\limits_{j=1}^{n} x_j m(B_j)$ is defined as the integral of $x(s)$ over $S$ with respect to the measure $m$; the integral will be denoted by $\int_S x(s) m(\dd s)$ 
or, in short, by $\int_{S} x(s)$ or simply by $\int x(s)$ when no confusion can be expected. A real- (or complex-) valued function $x(s)$ defined $m$-a.e. on $S$ is said to be $m$-integrable or, in short, integrable over $S$ if there exists a sequence $\{x_n(s)\}$ of finitely-valued integrable functions converging to $x(s)$ $m$-a.e. and such that 
\[
\lim\limits _{n, k \rightarrow \infty} \int_{S} | x _{n} (s)- x _{k} (s) | m(\dd s) = 0.
\]
\end{definitionu}

It is then proved that a finite $\lim\limits_{n\to\infty}\int_{S}x_{n}(s)m(\dd s)$ exists and the value of this limit is independent of the choice of the approximating sequence $\{x_{n}(s)\}$. The value of the integral $\int_{S}x(s)m(\dd s)$ over $S$ with respect to the measure $m$ is, by definition, given by $\lim\limits_{n\to\infty}\int_{S}x_{n}(s)m(\dd s)$. We shall sometimes abbreviate the notation $\int_{S}x(s)m(\dd s)$ to $\int x(s)m(\dd s)$ or to $\int x(s)$. 

\yossub{Properties of the Integral}

\begin{enumerate}[i)]
 \item If $x(s)$ and $y(s)$ are integrable, then $\alpha x(s) + \beta y(s)$ is integrable and $\int_{S} (\alpha x(s) + \beta y(s)) m(\dd s) = \alpha \int_{S} x(s) m(\dd s) + \beta \int_{S} y(s) m(\dd s).$ 
 \item $x(s)$ is integrable iff $|x(s)|$ is integrable. 
 \item If $x(s)$ is integrable and $x(s) \geqq 0$ a.e., then $\int_{S} x(s) m(\dd s) \geqq 0$, and the equality sign holds iff $x(s) = 0$ a.e. 
 \item If $x(s)$ is integrable, then the function $X(B) = \int_{B} x(s) m(\dd s)$ is $\sigma$-additive, that is, $X\left(\sum\limits_{j=1}^{\infty} B_j\right) = \sum\limits_{j=1}^{\infty} X(B_j)$ for any disjoint sequence $\{B_j\}$ of sets $\in \mathfrak{B}$. Here $\int_{B} x(s) m(\dd s) = \int_{S} C_B(s) x(s) m(\dd s)$, where $C_B(s)$ is the defining function of the set $B$, that is, 
 \[
 C _{B} (s) = 1 \text{ for } s \in B \text{ and } C _{B} (s) = 0 \text{ for } s \in S- B.
 \]
 \item $X(B)$ in iv) is absolutely continuous with respect to $m$ in the sense that $m(B) = 0$ implies $X(B) = 0$. This condition is equivalent to the condition that $\lim\limits_{m(B) \to 0} X(B) = 0$ uniformly in $B \in \mathfrak{B}$. 
\end{enumerate}

\begin{lemman}[The Lebesgue-Fatou Lemma] Let $\{x_{n}(s)\}$ be a sequence of real-valued integrable functions. If there exists a real-valued integrable function $x(s)$ such that $x(s) \geqq x_{n}(s)$ a.e. for $n = 1, 2, \ldots$ (or $x(s) \leqq x_{n}(s)$ a.e. for $n = 1, 2, \ldots$), then 
\begin{gather*}\int_{S} \left(\varlimsup\limits_{n \rightarrow \infty} x _{n} (s)\right) m(d s) \geqq \varlimsup\limits_{n \rightarrow \infty} \int_{S} x _{n} (s) m(d s) \\\left(\text{ or } \int_{S} \left(\varliminf\limits_{n \to \infty} x _{n} (s)\right) m(d s) \leqq \varliminf\limits_{n \to \infty} \int_{S} x _{n} (s) m(d s)\right), 
\end{gather*}

under the convention that if $\varlimsup\limits_{n\to\infty}x_n(s)$ (or $\varliminf\limits_{n\to\infty}x_n(s)$) is not integrable, we understand that $\varlimsup\limits_{n\to\infty}\int_{S}x_n(s)m(\dd s)=-\infty$ (or $\varliminf\limits_{n\to\infty}\int_{S}x_n(s)m(\dd s)=\infty$). 
\end{lemman}

\begin{definitionu} Let $(S, \mathfrak{B}, m)$ and $(S', \mathfrak{B}', m')$ be two measure spaces. We denote by $B \times B'$ the smallest $\sigma$-ring of subsets of $S \times S'$ which contains all the sets of the form $B \times B'$, where $B \in B$, $B' \in B'$. It is proved that there exists a uniquely determined $\sigma$-finite, $\sigma$-additive and non-negative measure $m \times m'$ defined on $B \times B'$ such that 
\[
(m \times m ^{\prime}) (B \times B ^{\prime}) = m(B) m ^{\prime} (B ^{\prime}).
\]
$m \times m'$ is called the product measure of $m$ and $m'$. We may define the $\mathfrak{B} \times \mathfrak{B}'$-measurable functions $x(s, s')$ defined on $S \times S'$, and the $m \times m'$-integrable functions $x(s, s')$. The value of the integral over $S \times S'$ of an $m \times m'$-integrable function $x(s, s')$ will be denoted by 
\[
\iint_ {S \times S ^{\prime}} x (s, s ^{\prime}) (m \times m ^{\prime}) (\dd s \dd s ^{\prime}) \quad \text{ or } \quad \iint_ {S \times S ^{\prime}} x (s, s ^{\prime}) m(\dd s) m ^{\prime} (\dd s ^{\prime}).
\]
\end{definitionu}

\begin{theoremn}[The Fubini-Tonelli Theorem] $A$ $\mathfrak{B} \times \mathfrak{B}'$-measurable function $x(s, s')$ is $m \times m'$-integrable over $S \times S'$ iff at least one of the iterated integrals 
\[
\int_ {S ^{\prime}} \left\{\int_{S} | x (s, s ^{\prime}) | m(\dd s) \right\} m ^{\prime} (\dd s ^{\prime}) \quad \text{ and } \quad \int_{S} \left\{\int_ {S ^{\prime}} | x (s, s ^{\prime}) | m ^{\prime} (\dd s ^{\prime}) \right\} m(\dd s)
\]

is finite; and in such a case we have 
\end{theoremn}
\[
\begin{aligned} \iint_ {S \times S ^{\prime}} x (s, s ^{\prime}) m(\dd s) m ^{\prime} (\dd s ^{\prime}) & = \int_ {S ^{\prime}} \left\{\int_{S} x (s, s ^{\prime}) m(\dd s) \right\} m ^{\prime} (\dd s ^{\prime}) \\ & = \int_{S} \left\{\int_ {S ^{\prime}} x (s, s ^{\prime}) m ^{\prime} (\dd s ^{\prime}) \right\} m(\dd s). \end{aligned}
\]

\yossub{Topological Measures}

\begin{definitionu} Let $S$ be a locally compact space, e.g., an $n$-dimensional Euclidean space $\mathbb{R}^{n}$ or a closed subset of $\mathbb{R}^{n}$. The Baire subsets of $S$ are the members of the smallest $\sigma$-ring of subsets of $S$ which contains every compact $G_{\delta}$-set, i.e., every compact set of $S$ which is the intersection of a countable number of open sets of $S$. The Borel subsets of $S$ are the members of the smallest $\sigma$-ring of subsets of $S$ which contains every compact set of $S$. 
\end{definitionu}

If $S$ is a closed subset of a Euclidean space $\mathbb{R}^{n}$, the Baire and the Borel subsets of $S$ coincide, because in $\mathbb{R}^{n}$ every compact (closed bounded) set is a $G_\delta$-set. If, in particular, $S$ is a real line $\mathbb{R}^{1}$ or a closed interval on $\mathbb{R}^{1}$, the Baire (= Borel) subsets of $S$ may also be defined as the members of the smallset $\sigma$-ring of subsets of $S$ which contains half open intervals $(a, b]$. 

\begin{definitionu} Let $S$ be a locally compact space. Then a non-negative Baire (Borel) measure on $S$ is a $\sigma$-additive measure defined for every Baire (Borel) subset of $S$ such that the measure of every compact set is finite. The Borel measure $m$ is called regular if for each Borel set $B$ we have 
\[
m(B) = \inf\limits _{U \supseteq B} m(U)
\]
where the infimum is taken over all open sets $U$ containing $B$. We may also define the regularity for Baire measures in a similar way, but it turns out that a Baire measure is always regular. It is also proved that each Baire measure has a uniquely determined extension to a regular Borel measure. Thus we shall discuss only Baire measures. 
\end{definitionu}

\begin{definitionu} $A$ complex-valued function $f(s)$ defined on a locally compact space $S$ is a Baire function on $S$ if $f^{-1}(B)$ is a Baire set of $S$ for every Baire set $B$ in the complex plane $\mathbb{C}^{1}$. Every continuous function is a Baire function if $S$ is a countable union of compact sets. A Baire function is measurable with respect to the $\sigma$-ring of all Baire sets of $S$. 
\end{definitionu}

\yossub{The Lebesgue Measure}

\begin{definitionu} Suppose $S$ is the real line $\mathbb{R}^{1}$ or a closed interval of $\mathbb{R}^{1}$. Let $F(x)$ be a monotone non-decreasing function on $S$ which is continuous from the right: $F(x) = \inf\limits_{x<y} F(y)$. Define a function $m$ on half closed intervals $(a, b]$ by $m((a, b]) = F(b)- F(a)$. This $m$ has a uniquely determined extension to a non-negative Baire measure on $S$. The extended measure $m$ is finite, i.e., $m(S) < \infty$ iff $F$ is bounded. If $m$ is the Baire measure induced by the function $F(s) = s$, then $m$ is called the Lebesgue measure. The Lebesgue measure in $\mathbb{R}^{n}$ is obtained from the $n$-tuple of the one-dimensional Lebesgue measures through the process of forming the product measure. 
\end{definitionu}

Concerning the Lebesgue measure and the corresponding Lebesgue integral, we have the following two important theorems: 

\begin{theorem}\label{thm:0.3.Theorem.1} Let $M$ be a Baire set in $\mathbb{R}^{n}$ whose Lebesgue measure $|M|$ is finite. Then, if we denote by $B \ominus C$ the symmetric difference of the set $B$ and $C: B \ominus C = B \cup C- B \cap C$, we have 
\[\lim\limits_{|h|\to 0}|(M + h)\ominus M| = 0, \text{ where } M + h = \{x\in \mathbb{R}^{n}; x = m + h,m\in M\}.\] 
Here $m + h = (m_{1} + h_{1},\dots ,m_{n} + h_{n})$ for $m = (m_1,\dots ,m_n)$ $h = (h_1,\dots ,h_n)$ and $|h| = \left(\sum\limits_{j=1}^{n}h_{j}\right)^{1/2}$.
\end{theorem}

\begin{theorem}\label{thm:0.3.Theorem.2} Let $G$ be an open set of $\mathbb{R}^{n}$. For any Lebesgue integrable function $f(x)$ in $G$ and $\varepsilon > 0$, there exists a continuous function $C_{\varepsilon}(x)$ in $G$ such that $\{x \in G; C_{\varepsilon}(x) \neq 0\}^a$ is a compact subset of $G$ and 
\[
\int_{G} | f (x)- C _{\varepsilon} (x) | \dd x < \varepsilon .
\]
\end{theorem}


\begin{remarku} Let $m$ be a Baire measure on a locally compact space $S$. A subset $Z$ of $S$ is called a set of $m$-measure zero if, for each $\varepsilon > 0$, there is a Baire set $B$ containing $Z$ with $m(B) < \varepsilon$. One can extend $m$ to the class of $m$-measurable sets, such a set being one which differs from a Baire set by a set of $m$-measure zero. Any property pertaining to a set of $m$-measure zero is said to hold $m$-almost everywhere ($m$-a.e.). One can also extend integrability to a function which coincides $m$-a.e. with a Baire function. 
\end{remarku}


\yossection{4}{Linear Spaces}
\label{sec:0.4}

\yossub{Linear Spaces}

\begin{definitionu} $A$ set $X$ is called a linear space over a field $K$ if the following conditions are satisfied: 
\begin{equation}\label{eq:0.4.1}
\parbox{0.90\linewidth}{$X$ is an abelian group (written additively),}
\end{equation}
A scalar multiplication is defined: to every element $x \in X$ and each $\alpha \in K$ there is associated an element of $X$, denoted by $\alpha x$, such that we have 
\begin{equation}\label{eq:0.4.2}
\begin{aligned}
&\alpha (x + y) = \alpha x + \alpha y (\alpha \in K; x, y \in X), (\alpha + \beta) x = \alpha x + \beta x (\alpha , \beta \in K; x \in X),
\\
&(\alpha \beta) x = \alpha (\beta x) (\alpha , \beta \in K; x \in X), 1 \cdot x = x (1 \text{ is the unit element of the field } K).
\end{aligned}
\end{equation}
\end{definitionu}

In the sequel we consider linear spaces only over the real number field $\mathbb{R}^{1}$ or the complex number field $\mathbb{C}^{1}$. A linear space will be said to be real or complex according as the field $K$ of coefficients is the real number field $\mathbb{R}^{1}$ or the complex number field $\mathbb{C}^{1}$. Thus, in what follows, we mean by a linear space a real or complex linear space. We shall denote by Greek letters the elements of the field of coefficients and by Roman letters the elements of $X$. The zero of $X$ (= the unit element of the additively written abelian group $X$) and the number zero will be denoted by the same letter $0$, since it does not cause inconvenience as $0 \cdot x = (\alpha- \alpha)x = \alpha x- \alpha x = 0$. The inverse element of the additively written abelian group $X$ will be denoted by $-x$; it is easy to see that $-x = (-1)x$. 

\begin{definitionu} The elements of a linear space $X$ are called vectors (of $X$). The vectors $x_{1}, x_{2}, \ldots, x_{n}$ of $X$ are said to be linearly independent if the equation $\sum\limits_{j=1}^{n} \alpha_{j} x_{j} = 0$ implies $\alpha_{1} = \alpha_{2} = \cdots = 0$. They are linearly dependent if such an equation holds where at least one coefficient is different from $0$. If $X$ contains $n$ linearly independent vectors, but every system of $(n + 1)$ vectors is linearly dependent, then $X$ is said to be of $n$-dimension. If the number of linearly independent vectors is not finite, then $X$ is said to be of infinite dimension. Any set of $n$ linearly independent vectors in an $n$-dimensional linear space constitutes a basis for $X$ 
and each vector $x$ of $X$ has a unique representation of the form $x = \sum\limits_{j=1}^{n} \alpha_j y_j$ in terms of the basis $y_1, y_2, \ldots, y_n$. A subset $M$ of a linear space $X$ is called a linear subspace or, in short, a subspace, if whenever $x, y \in M$, the linear combinations $\alpha x + \beta y$ also belong to $M$. $M$ is thus a linear space over the same coefficient field as $X$. 
\end{definitionu}

\yossub{Linear Operators and Linear Functionals}

\begin{definitionu} Let $X$, $Y$ be linear spaces over the same coefficient field $K$. A mapping $T: x \to y = T(x) = Tx$ defined on a linear subspace $D$ of $X$ and taking values in $Y$ is said to be linear, if 
\[
T \left(\alpha x _{1} + \beta x _{2}\right) = \alpha \left(T x _{1}\right) + \beta \left(T x _{2}\right).
\]
\end{definitionu}

The definition implies, in particular, 
\[
T \cdot 0 = 0, T (- x) =- (T x).
\]

We denote 
\[
D = D (T), \{y \in Y; y = T x, x \in D (T) \} = R (T), \{x \in D (T); T x = 0 \} = N (T)
\]
and call them the domain, the range and the null space of $T$, respectively. $T$ is called a linear operator or linear transformation on $D(T) \subseteq X$ into $Y$, or somewhat vaguely, a linear operator from $X$ into $Y$. If the range $R(T)$ is contained in the scalar field $K$, then $T$ is called a linear functional on $D(T)$. If a linear operator $T$ gives a one-to-one map of $D(T)$ onto $R(T)$, then the inverse map $T^{-1}$ gives a linear operator on $R(T)$ onto $D(T)$: 
\[
T ^{-1} T x = x \text{ for } x \in D (T) \quad \text{ and } \quad T T ^{-1} y = y \text{ for } y \in R (T).
\]
$T^{-1}$ is the inverse operator or, in short, the inverse of $T$. By virtue of $T(x_{1}- x_{2}) = Tx_{1}- Tx_{2}$, we have the following. 

\begin{propositionu} $A$ linear operator $T$ admits the inverse $T^{-1}$ iff $Tx = 0$ implies $x = 0$. 
\end{propositionu}

\begin{definitionu} Let $T_{1}$ and $T_{2}$ be linear operators with domains $D(T_{1})$ and $D(T_{2})$ both contained in a linear space $X$, and ranges $R(T_{1})$ and $R(T_{2})$ both contained in a linear space $Y$. Then $T_{1} = T_{2}$ iff $D(T_{1}) = D(T_{2})$ and $T_{1}x = T_{2}x$ for all $x \in D(T_{1}) = D(T_{2})$. If $D(T_{1}) \subseteq D(T_{2})$ and $T_{1}x = T_{2}x$ for all $x \in D(T_{1})$, then $T_{2}$ is called an extension of $T_{1}$, and $T_{1}$ a restriction of $T_{2}$; we shall then write $T_{1} \subseteq T_{2}$. 
\end{definitionu}

Convention. The value $T(x)$ of a linear functional $T$ at a point $x \in D(T)$ will sometimes be denoted by $\langle x, T \rangle$, i.e. 
\[
T (x) = \langle x, T \rangle .
\]

\yossub{Factor Spaces}

\begin{propositionu} Let $M$ be a linear subspace in a linear space $X$. We say that two vectors $x_{1}, x_{2} \in X$ are equivalent modulo $M$ if $(x_{1}- x_{2}) \in M$ and write this fact symbolically by $x_{1} \equiv x_{2} \pmod{M}$. Then we have: 

\begin{enumerate}[(i)]
\item $x \equiv x \pmod{M}$,
\item if $x_{1} \equiv x_{2} \pmod{M}$, then $x_{2} \equiv x_{1} \pmod{M}$,
\item if $x_{1} \equiv x_{2} \pmod{M}$ and $x_{2} \equiv x_{3} \pmod{M}$, then $x_{1} \equiv x_{3} \pmod{M}$.
\end{enumerate}
\end{propositionu}
\begin{proof}
 \begin{enumerate}[(i)]
\item is clear since $x- x = 0 \in M$.
\item If $(x_{1}- x_{2}) \in M$, then $(x_{2}- x_{1}) =-(x_{1}- x_{2}) \in M$.
\item If $(x_{1}- x_{2}) \in M$ and $(x_{2}- x_{3}) \in M$, then $(x_{1}- x_{3}) = (x_{1}- x_{2}) + (x_{2}- x_{3}) \in M$. \qedhere
\end{enumerate}
\end{proof}

We shall denote the set of all vectors $\in X$ equivalent modulo $M$ to a fixed vector $x$ by $\xi_{x}$. Then, in virtue of properties (ii) and (iii), all vectors in $\xi_{x}$ are mutually equivalent modulo $M$. $\xi_{x}$ is called a class of equivalent (modulo $M$) vectors, and each vector in $\xi_{x}$ is called a representative of the class $\xi_{x}$. Thus a class is completely determined by any one of its representatives, that is, $y \in \xi_{x}$ implies that $\xi_{y} = \xi_{x}$. Hence, two classes $\xi_{x}, \xi_{y}$ are either disjoint (when $y \in \xi_{x}$) or coincide (when $y \in \xi_{x}$). Thus the entire space $X$ decomposes into classes $\xi_{x}$ of mutually equivalent (modulo $M$) vectors. 

\begin{theoremu} We can consider the above introduced classes (modulo $M$) as vectors in a new linear space where the operation of addition of classes and the multiplication of a class by a scalar will be defined through 
\[
\xi_{x} + \xi_{y} = \xi_{x+y}, \alpha \xi_{x} = \xi_{\alpha x}.
\]
\end{theoremu}

\begin{proof}
 The above definitions do not depend upon the choice of representatives $x, y$ of the classes $\xi_x, \xi_y$ respectively. In fact, if $(x_1- x) \in M$, $(y_1- y) \in M$, then 
\[
\begin{gathered}
(x _{1} + y _{1})- (x + y) = (x _{1}- x) + (y _{1}- y) \in M, (\alpha x _{1}- \alpha x) = \alpha (x _{1}- x) \in M.
\end{gathered}
\]

We have thus proved $\xi_{x_{1}+y_{1}}=\xi_{x+y}$ and $\xi_{ax_{1}}=\xi_{ax}$, and the above definitions of the class addition and the scalar multiplication of the classes are justified. 
\end{proof}

\begin{definitionu} The linear space obtained in this way is called the factor space of $X$ modulo $M$ and is denoted by $X/M$. 
\end{definitionu}

\yossub{References}

Topological Space: P. \textsc{Alexandrov}-H. \textsc{Hopf} \cite{alexandrov6}, N. \textsc{Bourbaki} \cite{bourbaki26}, J. L. \textsc{Kelley} \cite{kelley146}. 

Measure Space: P. R. \textsc{Halmos} \cite{halmos92}, S. \textsc{Saks} \cite{saks280}. 
